\documentclass[11pt]{article}
\pdfinfoomitdate=1
\pdftrailerid{}
\pdfsuppressptexinfo=15
\usepackage[a4paper,margin=0.94in]{geometry}
\usepackage[T1]{fontenc}
\usepackage{lmodern}
\usepackage{amsmath,amssymb,amsthm,booktabs,microtype}
\usepackage[numbers,sort&compress]{natbib}
\usepackage{cleveref}

\newtheorem{theorem}{Theorem}[section]
\newtheorem{corollary}[theorem]{Corollary}
\newtheorem{proposition}[theorem]{Proposition}
\theoremstyle{remark}
\newtheorem{remark}[theorem]{Remark}

\title{A Parity-Free Near-Alias\\Actual Direct-Tail Envelope}
\author{Bin Wang}
\date{August 2026}

\begin{document}
\maketitle

\begin{abstract}
RH-352 and RH-353 control normalized actual direct coefficients on selected
lower-even and boundary orders.  We remove the parity and order restrictions.
On one bounded-phase physical clock let
\[
 N_k=2k-L_k,\qquad p_{\sigma,k,n}=\tau_{\sigma,n}-a_n,
\]
where $0\le L_k\le2k-2$.  Put $s=qR$, $t=q_*R$,
$x=(\beta R)^2$, $u=s/\sqrt{x}$, and $v=t/\sqrt{x}$.  The source-locked
all-order caps imply the complete direct Wiener-tail bound
\[
 \mathfrak W_k:=x^{-N_k/2}\sum_{n\ge N_k}|p_{\sigma,k,n}|R^n
 \le
 \frac{q^{-2}\lambda^{-2\eta_k}}{1-s}\rho_N^ku^{-L_k}
 +\frac{48}{1-t}\rho_T^kv^{-L_k},
\]
where
\[
 \rho_N=\lambda^2u^2=\frac{r_H^2\lambda^3}{4},
 \qquad \rho_T=v^2=\frac1\lambda.
\]
Thus, for every $L_k=o(k)$,
\[
 \limsup_{k\to\infty}\mathfrak W_k^{1/k}
 \le\rho_N<\frac{1419857}{1600000}<1.
\]
This single actual theorem includes every odd and even order above the cut:
in particular the critical order, the first-lower order whenever $L_k\ge2$,
the complete upper-alias band, and every later order.  Moreover,
\[
 x^{-k}\sum_{N_k\le n<4k}
 \frac{|p_{\sigma,k,n}|R^n}{n}
\]
decays with the same root ceiling; so does the full logarithmic tail from
$N_k$ onward.  Exact linear-depth frontiers are derived for both
normalizations.  The unnormalized separate noisy majorant has superunit root
$\lambda^2(qR)^2>9604/7225$, which is a method boundary and not actual
noncancellation.  Since $p=q-d$, no full-trace $E_{\rm off}$ conclusion
follows without head-defect control.  RH-241, RH-288, Gates A--E, and every
Riemann-hypothesis claim remain open.
\end{abstract}

\section{One direct type at every order}

Fix the physical constants
\begin{equation}\label{eq:constants}
 r_H=\frac{17}{20},\qquad q=\frac12,\qquad R=\frac75,
 \qquad \frac{28}{17}<\lambda<\frac{17}{10}.
\end{equation}
The strict lower and upper multiplier bounds are certified in RH-262 and
RH-336, respectively
\citep{WangBoundaryBudget2026,WangProjectorMass2026}.  Put
\begin{equation}\label{eq:scales}
 \beta=\frac1{r_H\sqrt\lambda},\qquad
 q_*=\frac1{r_H\lambda},\qquad
 x=(\beta R)^2>1.
\end{equation}
Let $\sigma=\sigma_k\downarrow0$ lie on a bounded-phase physical clock:
\begin{equation}\label{eq:clock}
 \eta_k=k-\frac{\log(1/\sigma_k)}{2\log\lambda},
 \qquad \sup_k|\eta_k|<\infty,
 \qquad \sigma_k^{-1}=\lambda^{2(k-\eta_k)}.
\end{equation}

RH-282 constructs the modulus-complete normal spectral complement and proves,
for all $n\ge2$,
\begin{equation}\label{eq:tau-cap}
 \tau_{\sigma,n}=\operatorname{Tr}C_\sigma^n,
 \qquad |\tau_{\sigma,n}|\le\sigma^{-1}q^{n-2}.
\end{equation}
RH-267 proves the deterministic numerator envelope
\begin{equation}\label{eq:a-cap}
 |a_n|<48q_*^n\qquad(n\ge2).
\end{equation}
The latter theorem treats odd and even orders explicitly.  RH-288 and RH-340
lock their difference as one actual direct coefficient, at every finite order,
\begin{equation}\label{eq:p-type}
 \boxed{p_{\sigma,k,n}:=\tau_{\sigma,n}-a_n
 =q_{\sigma,k,n}-d_{\sigma,k,n}.}
\end{equation}
No orbit, parity, or alias decomposition is used below
\citep{WangSpectralTail2026,WangUnifiedEnvelope2026,WangGluing2026,
WangSynchronization2026}.

Define four subunit geometric parameters
\begin{equation}\label{eq:geometric}
 s=qR=\frac7{10},\qquad
 t=q_*R=\frac{28}{17\lambda},\qquad
 u=\frac{s}{\sqrt x}=qr_H\sqrt\lambda,
 \qquad v=\frac{t}{\sqrt x}=\lambda^{-1/2}.
\end{equation}
Indeed $s,t,v<1$, while
\[
 u^2=q^2r_H^2\lambda
 <\frac{4913}{16000}<1.
\]
The two natural root rates are
\begin{equation}\label{eq:rates}
 \rho_N=\lambda^2u^2=\frac{r_H^2\lambda^3}{4},
 \qquad \rho_T=v^2=\frac1\lambda.
\end{equation}
They satisfy
\begin{equation}\label{eq:rate-certificates}
 \rho_T<\rho_N<\frac{1419857}{1600000}<1,
 \qquad \rho_T<\frac{17}{28}<1.
\end{equation}
The upper bounds are those of RH-352.  For the ordering, the lower multiplier
bound gives $r_H\lambda^2>2$, hence
$\rho_N/\rho_T=r_H^2\lambda^4/4>1$
\citep{WangModulusCancellation2026}.

\section{A complete bottom-normalized direct tail}

Let $L_k$ be any integer sequence satisfying
\begin{equation}\label{eq:depth-domain}
 0\le L_k\le2k-2,
 \qquad N_k=2k-L_k\ge2.
\end{equation}
Define the bottom-normalized coefficient $\ell^1$ tail
\begin{equation}\label{eq:W}
 \mathfrak W_k(L_k)
 =x^{-N_k/2}\sum_{n\ge N_k}|p_{\sigma,k,n}|R^n.
\end{equation}
Both source series converge because $s,t<1$.

\begin{theorem}[Parity-free actual direct-tail envelope]\label{thm:W}
For every depth sequence in \eqref{eq:depth-domain},
\begin{equation}\label{eq:W-bound}
 \boxed{
 \mathfrak W_k(L_k)
 \le
 \frac{q^{-2}\lambda^{-2\eta_k}}{1-s}
 \rho_N^ku^{-L_k}
 +\frac{48}{1-t}\rho_T^kv^{-L_k}.}
\end{equation}
Consequently, if $L_k=o(k)$, then
\begin{equation}\label{eq:W-root}
 \boxed{
 \limsup_{k\to\infty}\mathfrak W_k(L_k)^{1/k}
 \le\rho_N<\frac{1419857}{1600000}<1.}
\end{equation}
\end{theorem}

\begin{proof}
The triangle inequality in \eqref{eq:p-type} and \eqref{eq:tau-cap} give
\begin{align*}
 x^{-N_k/2}\sum_{n\ge N_k}|\tau_{\sigma,n}|R^n
 &\le \frac{q^{-2}\sigma^{-1}}{1-s}
       \left(\frac{s}{\sqrt x}\right)^{N_k}\\
 &=\frac{q^{-2}\lambda^{-2\eta_k}}{1-s}
   (\lambda^2u^2)^k u^{-L_k}.
\end{align*}
Likewise \eqref{eq:a-cap} gives
\begin{align*}
 x^{-N_k/2}\sum_{n\ge N_k}|a_n|R^n
 &<\frac{48}{1-t}
       \left(\frac{t}{\sqrt x}\right)^{N_k}\\
 &=\frac{48}{1-t}(v^2)^kv^{-L_k}.
\end{align*}
These are the two terms in \eqref{eq:W-bound}.  Bounded phase and
$L_k=o(k)$ make every factor outside $\rho_N^k$ and $\rho_T^k$
subexponential.  Apply the standard root bound for a sum and
\eqref{eq:rate-certificates}.
\end{proof}

The theorem is stronger than a finite near-alias-band estimate: it controls
every order from $N_k$ to infinity.  It is also parity-free.  The only
order-specific information is the moving lower cut.

\section{The complete near-alias logarithmic band}

Define the alias-clock logarithmic tail and band
\begin{align}
 \mathfrak T_k(L_k)
 &=x^{-k}\sum_{n\ge N_k}
   \frac{|p_{\sigma,k,n}|R^n}{n},
 \label{eq:T}\\
 \mathfrak B_k(L_k)
 &=x^{-k}\sum_{N_k\le n<4k}
   \frac{|p_{\sigma,k,n}|R^n}{n}.
 \label{eq:B}
\end{align}

\begin{corollary}[Complete near-alias band and full logarithmic tail]
\label{cor:band}
For every depth sequence in \eqref{eq:depth-domain},
\begin{equation}\label{eq:T-from-W}
 \mathfrak B_k(L_k)\le\mathfrak T_k(L_k)
 \le\frac{x^{-L_k/2}}{N_k}\mathfrak W_k(L_k).
\end{equation}
In particular, if $L_k=o(k)$, then
\begin{equation}\label{eq:B-root}
 \boxed{
 \limsup_{k\to\infty}\mathfrak T_k(L_k)^{1/k}
 \le\rho_N<1,\qquad
 \limsup_{k\to\infty}\mathfrak B_k(L_k)^{1/k}
 \le\rho_N<1.}
\end{equation}
\end{corollary}

\begin{proof}
The first inequality is nonnegativity.  Since $n\ge N_k$ and
$k-N_k/2=L_k/2$,
\[
 \mathfrak T_k(L_k)
 \le\frac{x^{-k}}{N_k}\sum_{n\ge N_k}|p_{\sigma,k,n}|R^n
 =\frac{x^{-L_k/2}}{N_k}\mathfrak W_k(L_k).
\]
For $L_k=o(k)$, the additional factors have $k$th root tending to one.
Apply \cref{thm:W}.
\end{proof}

The band in \eqref{eq:B} contains both parities, the critical order $2k$,
the first-lower order $2k-2$ whenever it lies above the lower cut, and every
upper-alias order $2k<n<4k$.  This statement concerns the direct coefficient
$p$ only.  It does not assert all-order decompositions of the form
$Y+\mathcal P-S$.

\section{Linear-depth frontiers}

The source estimates allow a strict extension beyond sublinear depth.  Let
\begin{equation}\label{eq:ell}
 \ell=\limsup_{k\to\infty}\frac{L_k}{k}\in[0,2].
\end{equation}

\begin{theorem}[Exact source-bound root frontiers]\label{thm:frontiers}
For the two normalized tails,
\begin{align}
 \limsup_{k\to\infty}\mathfrak W_k(L_k)^{1/k}
 &\le\max\{\rho_Nu^{-\ell},\rho_Tv^{-\ell}\},
 \label{eq:W-linear}\\
 \limsup_{k\to\infty}\mathfrak T_k(L_k)^{1/k}
 &\le\max\{\rho_Ns^{-\ell},\rho_Tt^{-\ell}\}.
 \label{eq:T-linear}
\end{align}
Define
\begin{equation}\label{eq:alphas}
 \alpha_{\rm nat}
 =\frac{\log(1/\rho_N)}{\log(1/u)},
 \qquad
 \alpha_{\rm alias}
 =\frac{\log(1/\rho_N)}{\log(1/s)}.
\end{equation}
Then $0<\alpha_{\rm nat}<\alpha_{\rm alias}<2$.
The bound \eqref{eq:W-linear} is strictly subunit for
$\ell<\alpha_{\rm nat}$, and \eqref{eq:T-linear} is strictly subunit for
$\ell<\alpha_{\rm alias}$.
\end{theorem}

\begin{proof}
Apply \eqref{eq:W-bound} along the definition of the upper limit.  This gives
\eqref{eq:W-linear}.  Combining \eqref{eq:T-from-W} with
$x^{-\ell/2}u^{-\ell}=s^{-\ell}$ and
$x^{-\ell/2}v^{-\ell}=t^{-\ell}$ gives \eqref{eq:T-linear}.

The noisy thresholds in \eqref{eq:alphas} are positive because
$\rho_N<1$, and $u<s<1$ gives
$\alpha_{\rm nat}<\alpha_{\rm alias}$.  At depth two,
\[
 \rho_Nu^{-2}=\lambda^2>1,
 \qquad
 \rho_Ns^{-2}=\frac{r_H^2\lambda^3}{R^2}
 >\frac{R}{r_H}=\frac{28}{17}>1,
\]
so both noisy thresholds are below two.  The target threshold for
\eqref{eq:W-linear} is exactly two because $\rho_T=v^2$.  For the alias
normalization, $\rho_Tt^{-2}=x^{-1}<1$, so the target term is still subunit
through the whole permitted interval $[0,2]$.  Hence the noisy terms determine
the stated frontiers.
\end{proof}

Substitution of the archived physical decimal for $\lambda$ gives the
diagnostics
\[
 \alpha_{\rm nat}=0.263953\ldots,
 \qquad \alpha_{\rm alias}=0.441576\ldots.
\]
They are descriptive evaluations of the exact formulas, not interval
certificates and not finite evidence for the theorem.

At the alias frontier the polynomial factor in the logarithmic tail still
matters.

\begin{proposition}[Critical alias-depth source convergence]
\label{prop:critical}
If
\begin{equation}\label{eq:critical-condition}
 \rho_N^ks^{-L_k}=o(N_k),
 \qquad \rho_T^kt^{-L_k}=o(N_k),
\end{equation}
then $\mathfrak T_k(L_k)\to0$ and hence
$\mathfrak B_k(L_k)\to0$.  In particular,
\begin{equation}\label{eq:critical-linear}
 L_k=\alpha_{\rm alias}k+O(1)
\end{equation}
implies $\mathfrak T_k(L_k)=O(k^{-1})$.
\end{proposition}

\begin{proof}
Insert \eqref{eq:W-bound} into \eqref{eq:T-from-W}; the two terms are fixed
bounded factors times the ratios in \eqref{eq:critical-condition}.  Under
\eqref{eq:critical-linear}, the noisy numerator is bounded because
$\rho_N=s^{\alpha_{\rm alias}}$.  The target numerator remains exponentially
small because its root at $\alpha_{\rm alias}<2$ is strictly subunit.
Finally $N_k\asymp k$.
\end{proof}

The critical statement is an upper-bound theorem only.  It does not assert
that the source estimate is sharp for the actual coefficient.

\section{Normalization and claim boundary}

Remove the factor $x^{-k}$ from the near-alias logarithmic budget.  At every
sublinear depth, the RH-282 separate noisy majorant has root
\begin{equation}\label{eq:raw-root}
 \lambda^2(qR)^2=x\rho_N.
\end{equation}
The certified lower multiplier bound gives
\begin{equation}\label{eq:raw-superunit}
 \lambda^2(qR)^2
 >\left(\frac{28}{17}\right)^2\left(\frac7{10}\right)^2
 =\frac{9604}{7225}>1.
\end{equation}
The separate deterministic majorant has root $(q_*R)^2<1$.  Therefore the
two absolute source caps cannot prove unnormalized direct-band closure.  This
is a strict scoped failure of the modulus-cap method, not a lower bound for
$p=\tau-a$: additional signed or complex cancellation may be present.

The normalized theorem also does not close the full-trace off-alias budget.
Equation \eqref{eq:p-type} gives only
\begin{equation}\label{eq:pqd}
 p_{\sigma,k,n}=q_{\sigma,k,n}-d_{\sigma,k,n}.
\end{equation}
RH-340 proves that transfer between direct and full-trace prefix budgets
requires a same-clock estimate for the head/counterloop defect $d$.  No such
estimate is supplied here.  Thus the new parity-free direct theorem is not a
theorem for $E_{\rm off}$
\citep{WangSynchronization2026,WangBoundaryGap2026}.

The executable artifact checks the exact rational identities in
\eqref{eq:geometric}--\eqref{eq:rate-certificates}, the two finite geometric
majorants, the linear root functions, and the superunit raw certificate.  Its
finite rows use the in-range rational fixture $\lambda=5/3$ and both odd and
even lower cuts.  They reproduce formulas only; they are not observations of
$\tau$, $a$, $p$, the noisy operator, or an asymptotic sequence.

RH-354 proves an actual normalized direct-coefficient tail theorem.  It does
not control the low prefix $2\le n<N_k$, prove head transport, identify the
direct budget with full $E_{\rm off}$, close the RH-241 moving noisy all-order
envelope or coefficient bridge, or activate RH-288.  Gates A--E remain
false/open.  This paper constructs no Hilbert--Polya operator, identifies no
Riemann zero, proves no von Mangoldt-weighted prime-power trace, proves no
completed-zeta divisor equality, and does not prove the Riemann Hypothesis.

\bibliographystyle{plainnat}
\bibliography{references}

\end{document}
